\documentclass[journal]{IEEEtran}
\usepackage[T1]{fontenc}
\usepackage[utf8]{inputenc}
\usepackage{amsmath,amssymb,booktabs}
\usepackage{graphicx}
\title{Physics-Constrained Bayesian Reconstruction of Unobserved Power-System States from Sparse PMUs}
\author{Working draft -- Paper 2 of the sparse-PMU research program}

\begin{document}
\maketitle

\begin{abstract}
Sparse synchrophasor coverage turns transmission-system state estimation into
a functional inverse problem: the individual operating-point perturbations
that explain the measurements may not be unique, while the unobserved voltage
field can remain nearly determined. We study this distinction on a validated
PowerDynamics IEEE-39 benchmark with eight observed PMUs and 31 hidden buses.
The estimator combines the nominal dynamic model with a constrained nonlinear
AC maximum-a-posteriori (MAP) equilibrium solve. In the corrected M6
operating-point campaign, the nuisance vector has 43 coordinates and a median
local rank of 25, yet the hidden-voltage static closure is 0.974, 0.981, and
0.983 for mismatch scales 0.5, 1.0, and 1.5, respectively. A subsequent causal
streaming test shows that directly low-pass filtering dynamic PMU measurements
does not provide a safe online center: online closure is negative, convergence
is 0.8125, and nominal TVE increases materially. The results support
functional, rather than parameter, identifiability as the central criterion
for sparse-PMU reconstruction, while leaving causal center extraction and
joint uncertainty calibration as open problems.
\end{abstract}

\begin{IEEEkeywords}
Power-system state estimation, sparse PMUs, functional identifiability,
nonlinear AC MAP, operating-point mismatch, Bayesian estimation.
\end{IEEEkeywords}

\section{Introduction}
Eight PMUs observe buses $\mathcal P=\{2,5,6,10,19,22,29,39\}$ of the
39-bus system. The target is the hidden voltage field $V_{\mathcal U}$ on the
31 unobserved buses. This paper separates two questions that are often
conflated: whether nuisance injections are uniquely identifiable and whether
the hidden electrical field is identifiable as a functional of the data.

The paper makes four contributions. First, it formulates sparse-PMU
reconstruction with explicit AC constraints and a physical nuisance basis.
Second, it connects constrained local information diagnostics to hidden-voltage
functionals. Third, it reports a reproducible mismatch campaign on a
PowerDynamics-native IEEE-39 plant. Fourth, it documents the boundary between
an effective static MAP recentering and an unsafe raw causal low-pass extension.

\section{Problem Formulation}
Let $y_{\mathcal P}=h_{\mathcal P}(V)+v$ denote voltage and terminal-current
phasors at the eight PMUs. The operating-point nuisance vector is
\[
d=[\Delta P_{PQ},\Delta Q_{PQ},\Delta P_{PV}],\qquad \dim d=43,
\]
with bus 31 as the slack and bus 39 as a PV bus. The estimator solves
\begin{align}
\min_{V,d}\;&\tfrac12\|y_{\mathcal P}-h_{\mathcal P}(V)\|_{R^{-1}}^2
 +\tfrac12\|d-d_0\|_{\Gamma_d^{-1}}^2\\
\text{s.t. }&g_{AC}(V,d)=0,\quad \theta_{31}=0.
\end{align}
The hidden output is $V_{\mathcal U}$, not the nuisance vector itself.

\section{Functional Identifiability}
The constrained local map from nuisance coordinates to PMU outputs can have
rank $r<43$. E06-H gives median $r=25$ and nullspace dimension 18. The relevant
criterion is instead whether compatible nuisance perturbations are annihilated
by the hidden-voltage functional, i.e., whether the constrained nullspace is
approximately contained in the kernel of the hidden-output differential.
This permits non-unique physical explanations with a stable reconstructed
electrical field.

\section{Validated Benchmark and Estimator}
The benchmark contains 114 differential and 78 algebraic states, uses the
PowerDynamics IEEE-39 network primitives, and retains exact analytic AC and PMU
Jacobians. Static parity, equilibrium, linearization, and measurement-operator
audits precede the mismatch experiments. The physical nuisance basis keeps PV
reactive power as a solved quantity and treats the slack explicitly.

\section{Experiments}
\subsection{Static M6 gate (E06-H)}
Fresh DEV and TEST plants use the corrected basis, exact nominal network
admittance, and a nonlinear AC MAP. The frozen TEST split contains 60 cases,
20 per nonzero mismatch scale. The nominal center, corrected MAP, oracle
reference, AC residual, convergence, rank, and covariance are reported in the
E06-H result files.

\subsection{Causal streaming gate (E06-I)}
Fresh streaming splits contain 40 DEV and 120 TEST trajectories, including a
nominal $m=0$ control. The DEV-only grid compares trailing means, a robust
Huber trailing mean, and an exponential moving average over 10, 30, and 60
frames, with MAP updates every 1, 3, 10, or 30 frames. The selected causal
configuration is evaluated without future samples or hidden truth.

\section{Results}
\begin{table}[t]
\centering
\caption{Corrected static M6 reconstruction and causal streaming outcome.}
\label{tab:summary}
\begin{tabular}{lccc}
\toprule
Mismatch & Nominal TVE & E06-H MAP TVE & E06-H closure\\
\midrule
$m=0.5$ & 0.19885\% & 0.00526\% & 0.974\\
$m=1.0$ & 0.38834\% & 0.00749\% & 0.981\\
$m=1.5$ & 0.57795\% & 0.00974\% & 0.983\\
\bottomrule
\end{tabular}
\end{table}

E06-H yields strong static recovery despite nuisance non-identifiability. The
causal E06-I extension selected a 60-frame Huber window and 30-frame update
cadence, but obtained negative online closure on all nonzero scales. The
nominal control exposed the failure: R2 TVE rose to approximately 5.69\%, while
the frozen B2 reference remained below 0.02\%. The failure is therefore
consistent with dynamic-deviation leakage into the center target, not with a
failure of the static AC MAP itself.

\section{Uncertainty and Computational Cost}
The current Laplace approximation propagates static-center and Kalman
deviation covariance while omitting their cross-covariance. Marginal coverage
is not sufficient for calibration; the reported E06-H and E06-I NEES-like
diagnostics remain poor. Static MAP solves take approximately 61.7 ms in E06-H;
the selected E06-I update takes 38.3 ms median and 82.2 ms at p95. The latter is
amortized to 2.31 ms per frame at a 30-frame cadence, but each update remains a
blocking operation and is not a hard-real-time claim.

\section{Discussion}
The evidence supports a hierarchy in which a nonlinear static AC MAP supplies
the slow electrical center and a linear estimator supplies fast deviations.
The E06-I result shows that the center cannot be obtained by low-pass filtering
raw dynamic PMU channels. A causal innovation/deviation-compensated extractor
is required before claiming an online adaptive estimator. Parameter
identifiability is not the correct success criterion when the hidden voltage
functional is stable across the nuisance nullspace.

\section{Limitations and Next Experiment}
The study uses one network and one PMU placement. Uncertainty is not yet
calibrated, and the causal center problem remains open. The next experiment is
a preregistered M6-only innovation-based center extractor, followed by a
frozen transfer check to M1 and M7 only if the M6 nominal-safety gate passes.

\section{Conclusion}
Sparse PMUs can support physically constrained reconstruction of an unobserved
voltage field under operating-point mismatch even when the nuisance injections
are non-identifiable. The corrected static AC MAP recovers approximately 98\%
of the M6 mismatch error. The causal failure is equally informative: raw
low-pass centering is unsafe, so online adaptation must explicitly separate
slow operating-point structure from fast dynamic deviations.

\section*{Data and Code Availability}
Manifests, scripts, tests, and result tables are stored in the repository paths
listed in the accompanying claim ledger. The generated physical trajectories
remain evaluation artifacts; estimator inputs exclude hidden truth.

\section*{Author Contributions, Funding, Conflicts, and AI Disclosure}
To be completed before submission using the venue-specific declarations. No
human or animal subjects are involved.

\bibliographystyle{IEEEtran}
\bibliography{references}
\end{document}
